\begin{env}[7.9.7]
\label{IV.2.s7.9.7-fr}
Consideremos, por una parte, una subcategoría plena $\boldsymbol{C}$ de la categoría
de los preesquemas localmente noetherianos y, por otra, una propiedad
$\boldsymbol R(A)$, sujetas a las condiciones siguientes:

$1^\circ$ Para todo $X\in\boldsymbol{C}$, todo preesquema localmente de tipo finito
sobre $X$ pertenece a $\boldsymbol{C}$. Para todo anillo noetheriano $A$ tal que
$\operatorname{Spec}(A)\in\boldsymbol{C}$ y toda parte multiplicativa $S$ de
$A$, se tiene $\operatorname{Spec}(S^{-1}A)\in\boldsymbol{C}$.

$2^\circ$ Para todo $X\in\boldsymbol{C}$, el conjunto
$U_{\mathbf{R}}(X)$ de los $x\in X$ tales que
$\boldsymbol R(\mathcal O_{X,x})$ sea verdadera es abierto en $X$.

$3^\circ$ Para todo anillo local noetheriano $A$ tal que
$\operatorname{Spec}(A)\in\boldsymbol{C}$ y todo elemento regular $t$ del ideal
maximal de $A$, $\boldsymbol R(A/tA)$ implica $\boldsymbol R(A)$.

Denotemos entonces, como en (7.5.0), por $\boldsymbol P(Z,k)$, para un cuerpo $k$
y un $k$-preesquema $Z\in\boldsymbol{C}$, la propiedad siguiente:

\oldpage[IV-2]{221}

«Para toda extensión \emph{finita} $k'$ de $k$, todos los anillos locales de
$Z\otimes_k k'$ satisfacen la propiedad $\boldsymbol R$.»

Supondremos además que $\boldsymbol R$ satisface la condición siguiente:

$4^\circ$ Si $\boldsymbol P(Z,k)$ es verdadera, entonces, para toda extensión
\emph{de tipo finito} $k''$ de $k$, todos los anillos locales de
$Z\otimes_k k''$ satisfacen la propiedad $\boldsymbol R$.

Se observará que estas condiciones se satisfacen si se toma como
$\boldsymbol{C}$ la categoría de los preesquemas \emph{excelentes} y como
$\boldsymbol R$ una de las propiedades (i) a (viii) de (7.5.3); para la condición
$1^\circ$, esto se deduce de (7.8.3, (ii)), y para las condiciones $2^\circ$ y
$3^\circ$, de los razonamientos de (7.5.3), teniendo en cuenta que todo anillo
excelente es catenario. Por último, la condición $4^\circ$ se deduce,
para (i), (ii) y (iii), de (6.7.1); para (iv), (v) y (vi), de (6.7.7), y para
(viii), de (4.6.1); en cuanto a la propiedad (vii) de (7.5.3),
la propiedad correspondiente $\boldsymbol P(Z,k)$ implica que $Z$ es localmente
íntegro (por ser localmente noetheriano) y que cada uno de los subpreesquemas
íntegros cuya suma es $Z$ es geométricamente íntegro, en virtud de (4.5.9)
y (4.6.1); de ello se deduce también la propiedad $4^\circ$ anterior en este caso.

Con estas notaciones e hipótesis:
\end{env}

